\chapter{Results and Discussion}
\label{ch:results}

The completed application successfully implements a browser-playable
Minesweeper game with persistent accounts, a global leaderboard, personal
statistics, and game history, on top of a game engine that runs entirely on
the client. In its current form, the system:

\begin{itemize}[nosep]
    \item Correctly guarantees a safe first click and its neighbourhood on
    every new game, regardless of difficulty.
    \item Correctly cascades reveals through zero-adjacency regions using an
    iterative, non-recursive flood-fill, avoiding recursion-depth concerns
    on the larger $16\times16$ Advanced board.
    \item Supports chording and a single-use shield ability as additional
    layers of player strategy beyond classic reveal/flag play.
    \item Enforces a strict separation between client-side gameplay and
    server-side score authority: the score shown to the player during play
    is never the value trusted for the leaderboard, which is instead
    recomputed and bounds-checked independently by \texttt{save-score.php}.
    \item Presents personal statistics, a full game history, and a global
    leaderboard, each filterable by difficulty, giving the player both
    short-term feedback and long-term progress tracking.
\end{itemize}

The decision to implement the game engine as a set of pure, testable
functions (\texttt{minesweeper.js}), separate from the React state
management in \texttt{useMinesweeper.js}, kept the core algorithms
(first-click safety, flood-fill, chording, scoring) independent of the UI
framework and straightforward to reason about and validate directly, as
reflected in the test cases in Chapter~\ref{ch:testing}. Likewise, treating
the client-computed score as untrusted input and recomputing it
server-side, rather than simply persisting whatever the client reports,
directly addresses the score-integrity problem identified in the
Introduction (Chapter~\ref{ch:introduction}).

The use of a single-page React frontend against a lightweight,
framework-free procedural PHP API kept the backend small and easy to audit,
while still demonstrating a complete client-server-database pipeline with
real authentication, authorization, and data-integrity concerns, which
suited the goals of the project.
